\documentclass{leveldoc}
\DocUses{foundations,slingshot,thomson_crt,first_coil,stray_field,earths_field}
\DocFingerprint{2da053221ec72703}

\begin{document}

\section{What the level asks}\label{sec:asks}

The game's text:
\leveldescription

This is the cathode-ray tube of \levelref{thomson_crt}, with its reference charge
$Q = \val{QB}$ at $(\val{xB}, \val{yB})$ now built in, in a world with $c = \val{c}$. Two electrons
($q = \val{q}$, $m = \val{m}$) are launched from $(\val{xZero}, \val{yZero})$ along $+x$: the slow
one, $T = \val{TSlow}$, must enter the one-cell spot $\val{xSSlow} \le x \le \val{xSHiSlow}$,
$\val{ySLoSlow} \le y \le \val{ySHiSlow}$; the fast one, $T = \val{TFast}$, the spot
$\val{xSFast} \le x \le \val{xSHiFast}$, $\val{ySLoFast} \le y \le \val{ySHiFast}$. The tube stands
in a uniform stray field $B_z = \val{B}$ (``\val{distName}'', \physref{2.3}). The player adds up to
three charges from the listed magnitudes, either sign, at nodes with
$\val{xRegLo} \le x \le \val{xRegHi}$, $\val{yRegLo} \le y \le \val{yRegHi}$. A flight ends when the
electron enters its spot, touches a charge (radius $\val{R}$), leaves the world, or reaches
$t = \val{tMax}$ (\xsecref{foundations}{sec:board}).

The field bends both electrons up, past their spots; the reference adds $Q = \val{QR}$ at
$(\val{xQ}, \val{yQ})$, which attracts them back down by about as much. The document marches the
four flights (with and without the added charge) and measures the field's shift and the
correction beam by beam, explaining why one charge can cancel the field for both energies only
approximately.

\section{Setup}\label{sec:setup}

The level file, field by field (the appendix reproduces it):
\setuptable

By \xeqref{foundations}{eq:launch} the slow electron has $\gamma_0 = \val{gammaSlow}$, momentum
$p = \val{pSlow}$ and speed $\val{vSlow}$; the fast one $\gamma_0 = \val{gammaFast}$,
$p = \val{pFast}$, speed $\val{vFast}$. In the stray field alone they would circle with the radii
$R_B = p/(\abs{q}B) = \val{RBSlow}$ and $\val{RBFast}$, by
\xeqref{first_coil}{eq:larmor-radius}.

\begin{figure}[htbp]
\centering
\begin{tikzpicture}
\begin{axis}[axis equal image, width=0.92\linewidth, xmin=-1.6, xmax=31.6, ymin=-1.6, ymax=21.6,
  xtick={0,5,...,30}, ytick={0,5,...,20}, tick label style={font=\small},
  xlabel={$x$}, ylabel={$y$}, label style={font=\small}]
  \draw[black!8] (0,0) grid (30,20);
  \draw[black!45, dashed] (-1,-1) rectangle (\rawval{xWorld},\rawval{yWorld});
  \draw[black!30, dotted] (\rawval{xRegLo},\rawval{yRegLo}) rectangle (\rawval{xRegHi},\rawval{yRegHi});
  \fill[green!50!black, fill opacity=0.25] (\rawval{xSSlow},\rawval{ySLoSlow}) rectangle (\rawval{xSHiSlow},\rawval{ySHiSlow});
  \fill[green!50!black, fill opacity=0.25] (\rawval{xSFast},\rawval{ySLoFast}) rectangle (\rawval{xSHiFast},\rawval{ySHiFast});
  \addplot[thin, dashed, red!70!black] table {builtslow.dat};
  \addplot[thin, dashed, blue!70!black] table {builtfast.dat};
  \addplot[thick, red!70!black] table {pathslow.dat};
  \addplot[thick, blue!70!black] table {pathfast.dat};
  \fill[blue!60!black] (\rawval{xB},\rawval{yB}) circle[radius=\rawval{R}];
  \fill[red!75!black] (\rawval{xQ},\rawval{yQ}) circle[radius=\rawval{R}];
  \fill (\rawval{xZero},\rawval{yZero}) circle[radius=0.15];
\end{axis}
\end{tikzpicture}
\caption{The built-in charge (blue), the reference's added charge (red), the spots (green) and
the flights in the field: slow (red) and fast (blue), with the added charge (solid) and without it
(dashed, both leaving the world).}
\label{fig:board}
\end{figure}

\section{Principles and model}\label{sec:model}

\begin{principles}
  \item The relativistic equation of motion \xeqref{foundations}{eq:eom} with the charges'
    Coulomb fields \xeqref{foundations}{eq:coulomb} and the uniform $B_z$ (\physref{2.3}).
  \item The energy in a static field, \xeqref{foundations}{eq:energy}; Liénard's power
    \xeqref{foundations}{eq:larmor}.
  \item The small-kick angle of a passing charge, \xeqref{thomson_crt}{eq:impulse}, and the circle
    in a uniform field, \xeqref{first_coil}{eq:larmor-radius}.
  \item The board's rules; the classical Runge--Kutta method, cubic Hermite interpolation and the
    error estimate from the step $2h$ (\xsecref{foundations}{sec:board},
    \xsecref{foundations}{sec:numerics}), as in \xsecref{slingshot}{sec:method}.
\end{principles}

Each electron starts at $(\val{xZero}, \val{yZero})$ with $\vect{p} = (p, 0)$ and moves by
\xeqref{foundations}{eq:eom} with $\vect{F} = q(\vect{E}_Q + \vect{v}\times B_z\uvec{z})$ until the
first of the board's events. The field is static, so
\begin{equation}\label{eq:energy-crt}
  W = \gamma mc^2 + \sum_Q \frac{qQ}{\abs{\vect{x} - \vect{r}_Q}}
\end{equation}
is conserved (the magnetic force does no work).

\begin{assumption}{The electrons' radiation}\label{as:radiation}
The electrons radiate the power \xeqref{foundations}{eq:larmor}, which the game, like this
document, neglects; along the reference flights it is $\val{radSlow}$ and $\val{radFast}$, fractions
$\val{radFracSlow}$ and $\val{radFracFast}$ of their kinetic energies.
\end{assumption}

\begin{assumption}{Fixed charges, point electrons, a uniform stray field}\label{as:fixed}
The charges stay put, the electrons are points and do not act on each other (each is flown
alone), and the stray field is uniform and static, an exact solution of Maxwell's equations
(\physref{2.3}). These are the game's model.
\end{assumption}

\section{Derivation}\label{sec:derivation}

\subsection{The march}\label{sec:march}

There is no closed form. Each flight is marched by \xeqref{foundations}{eq:rk4} at $h = \val{H}$
with Liénard's energy as one more state; the events are located on the cubic Hermite interpolant
by bisection, the error is estimated against the step $2h$, and \eqref{eq:energy-crt} checks
every step (\xsecref{stray_field}{sec:march}).

\subsection{The field's shift and the correction}\label{sec:shifts}

How far the stray field moves each beam, and how far the added charge moves it back, show on the
spots' line $x = \val{xSSlow}$. Marching each electron with the built-in charge only (without and
with the field), with the added charge without the field, and with both:
\begin{center}\small\val{lines}\end{center}
The columns: the crossing height without field (the tube of \levelref{thomson_crt}, on target);
in the field; the field's shift; the added charge's shift without the field; and the crossing with
both. A circle of radius $R_B$ departs from its tangent by about $L^2/2R_B$ over $L$ cells: here
$\val{sagSlow}$ and $\val{sagFast}$, close to the field's measured shifts $\val{shiftBSlow}$ and
$\val{shiftBFast}$ (the built-in charge has already bent the rays, which changes the path the
field acts along).

The two shifts scale differently with the energy. The magnetic bend goes as $1/p$, the inverse
radius $\abs{q}B/p$; the added charge's kick turns a ray by $\kappa/(Tb)$,
\xeqref{thomson_crt}{eq:impulse}, as $1/T \approx 2/(pv)$. From the slow beam to the fast one $p$
grows by the factor $\val{pRatio}$ but $pv$ by $\val{pvRatio}$ (nearly $T$'s $\val{TRatio}$), so a charge that
cancels the field for one energy over- or under-corrects the other. The reference balances the
two: it leaves the slow beam $\val{sumSlow}$ off its no-field crossing and the fast one
$\val{sumFast}$, within the one-cell spots' tolerance (the slow beam crosses the line just below its
spot and enters it through the bottom face as it climbs).

\section{The reference solution by hand}\label{sec:hand}

\begin{procedure}
  \item\label{st:slow} \emph{The slow electron.} The march of \S\ref{sec:march}: it enters its spot
    through the \val{faceSlow} face at $(\val{xSlow}, \val{ySlow})$ at $t = \val{tSlow}$, with
    $\vect{p} = (\val{pxSlow}, \val{pySlow})$ (step estimate $\val{errSlow}$; the energy drifts by
    $\val{wDevSlow}$).
  \item\label{st:fast} \emph{The fast electron.} Through the \val{faceFast} face at
    $(\val{xFast}, \val{yFast})$ at $t = \val{tFast}$, with $\vect{p} = (\val{pxFast}, \val{pyFast})$
    (step estimate $\val{errFast}$; energy drift $\val{wDevFast}$).
\end{procedure}

\begin{result}[The reference flights]
Both electrons enter their spots in the stray field. The level is solved.
\end{result}

\section{Comparison with the game}\label{sec:comparison}

\comparisontable

All trajectory differences are at most $\val{largest}$, the game's verified accuracy
(Foundations §6), against tolerances of ten times the larger of the two computations' estimates.
The radiated energies agree to about $10^{-9}$ relative; the game integrates Liénard's power by
three-point Gauss--Legendre on each step, we as a state of the march.

\section{Other solutions}\label{sec:others}

The tube as built fails in the field (the rows ``as built''): the slow electron leaves the world
through its \val{builtFaceSlow} edge at $x = \val{builtXSlow}$, the fast one through the
\val{builtFaceFast} edge at $y = \val{builtYFast}$. A correction must pull both beams down by about
the field's shifts, more for the slow beam than for the fast one; a positive charge below the rays
near the gun does so, as \S\ref{sec:shifts} shows, and the $1/T$ scaling of its kick against the
field's $1/p$ sets how well one charge can serve both. A map of all solutions would take the
game's verdicts, which this document does not reproduce.

\section{What the game computes}\label{sec:engine}

The game adds the stray field to the charges' fields (\physref{2.3}) and integrates the equations
of motion with its adaptive eighth-order Runge--Kutta method (\physref{3}, \physref{5},
\physref{6}, \physref{7}), each shot verified on its own. The problem is the same as here, and the
differences in the comparison are the two integrations' errors.

\section{Further reading}\label{sec:reading}

The tube itself: \levelref{thomson_crt}. Motion in a uniform magnetic field: \Jackson{§12.2}. A
stray field handled by a fold instead: \levelref{earths_field}.

\levelfileappendix

\end{document}
